\documentclass[../DoAn.tex]{subfiles}
\begin{document}

Chương 2 đã trình bày nền tảng lý thuyết về sinh trắc học hành vi và phân tích bốn công trình liên quan, qua đó chỉ ra khoảng trống nghiên cứu về một hệ thống tích hợp đồng thời xác thực đa phương thức, khai thác ngữ cảnh hoạt động, nhận dạng tập mở và cơ chế xác thực dự phòng. Chương 3 trình bày phương pháp đề xuất nhằm lấp khoảng trống đó, bắt đầu từ thiết kế tổng thể của hệ thống và quá trình xây dựng bộ dữ liệu phục vụ huấn luyện. Phần \ref{sec:tongquan_giaiphap} giới thiệu tổng quan pipeline của hệ thống và vai trò của từng thành phần. Phần \ref{sec:thietke} mô tả thiết kế thu thập dữ liệu trên nhiều người dùng và nhiều dòng thiết bị. Phần \ref{sec:appthuthu} trình bày ứng dụng BioAuth Data Collection được xây dựng để thu thập dữ liệu một cách nhất quán. Phần \ref{sec:tienxuly} mô tả quy trình tiền xử lý dữ liệu quán tính làm đầu vào cho giai đoạn huấn luyện mô hình.

\section{Tổng quan giải pháp}
\label{sec:tongquan_giaiphap}

Mục này mô tả tổng quan pipeline của hệ thống đề xuất nhằm giúp người đọc nắm được vị trí và vai trò của từng thành phần trước khi đi vào chi tiết ở các phần sau. Hình~\ref{fig:pipeline_overview} minh họa toàn bộ luồng xử lý từ thu thập dữ liệu đến triển khai trên thiết bị Android, được chia thành bốn giai đoạn nối tiếp nhau.

Hai giai đoạn đầu chuẩn bị dữ liệu cho mô hình. Giai đoạn thứ nhất là thu thập dữ liệu thông qua ứng dụng BioAuth Data Collection, tạo ra bộ dữ liệu thô gồm tín hiệu cảm biến quán tính chín kênh và các sự kiện tương tác màn hình từ 26 người tham gia trên nhiều dòng thiết bị khác nhau; mỗi mẫu chuyển động được gán nhãn theo một trong ba hoạt động đi bộ, đứng và ngồi. Giai đoạn thứ hai là tiền xử lý và trích xuất đặc trưng: tín hiệu quán tính được phân đoạn thành các cửa sổ cố định 200 mẫu (tương ứng khoảng 4 giây ở 50 Hz), còn dữ liệu chạm được chuyển hóa thành các vector đặc trưng động học mô tả thao tác chạm, cuộn và gõ phím.

Hai giai đoạn sau xây dựng và đưa mô hình vào hoạt động. Giai đoạn thứ ba là huấn luyện và hình thành hàm quyết định. Ba kiến trúc backbone gồm CNN 1D, CNN-LSTM và CNN-BiLSTM được so sánh trên nhánh quán tính, và được huấn luyện, đánh giá theo hai cấu hình ngữ cảnh hoạt động là chỉ đi bộ và toàn bộ ba hoạt động nhằm khảo sát ảnh hưởng của ngữ cảnh. Mỗi backbone được huấn luyện như bài toán phân loại nhiều người dùng chỉ nhằm học một không gian embedding, sau đó tầng phân loại được loại bỏ và chỉ giữ lại bộ mã hóa. Ở bước xác thực, mỗi chủ sở hữu được đăng ký bằng một tập nhỏ embedding đại diện (anchor) và điểm tin cậy được tính theo độ tương đồng cosine giữa embedding đầu vào với các anchor này, vận hành theo kịch bản nhận dạng tập mở. Nhánh touch sử dụng bộ phân loại Random Forest và được dung hợp với nhánh quán tính ở mức điểm số; cấu hình đa phương thức này giữ vai trò đối chứng, và như phân tích ở Chương 5, do không cải thiện hiệu năng trong kịch bản tập mở nên cấu hình triển khai cuối cùng chỉ sử dụng nhánh quán tính. Điểm tin cậy được làm mịn theo thời gian bằng cơ chế EWMA trước khi so với ngưỡng quyết định. Giai đoạn thứ tư là triển khai trên thiết bị: mô hình được tối ưu và chuyển đổi sang định dạng TensorFlow Lite rồi tích hợp vào ứng dụng BioAuth Authenticator để thực hiện xác thực liên tục trong nền, đồng thời kích hoạt cơ chế xác thực dự phòng bằng cử chỉ lắc khi điểm tin cậy giảm dưới ngưỡng.

\begin{figure}[H]
    \centering
    \includegraphics[width=\textwidth]{Hinh_ve/kientruc_tongquan.png}
    \caption{Tổng quan pipeline hệ thống xác thực hành vi đa phương thức}
    \label{fig:pipeline_overview}
\end{figure}

Như vậy, pipeline đề xuất kết nối bốn giai đoạn thành một vòng khép kín từ dữ liệu hành vi đến quyết định xác thực và phản ứng khi phát hiện bất thường. Các phần tiếp theo trình bày chi tiết từng giai đoạn, bắt đầu từ thiết kế thu thập dữ liệu trong Phần~\ref{sec:thietke}.

\section{Thiết kế thu thập dữ liệu}
\label{sec:thietke}

Bộ dữ liệu được xây dựng từ 26 người tham gia tình nguyện với độ tuổi và thói quen sử dụng điện thoại đa dạng. Mỗi người thực hiện thu thập trong nhiều phiên độc lập, mỗi phiên bao gồm ba hoạt động chuẩn hóa là đi bộ, đứng yên và ngồi với thời lượng mục tiêu mặc định 360 giây cho mỗi hoạt động (có thể điều chỉnh trong giao diện), tiếp theo là một đến hai lượt điền biểu mẫu tương tác màn hình chuẩn hóa.

Điểm khác biệt quan trọng trong thiết kế thu thập của đề tài là không áp đặt điều kiện tư thế hay môi trường cố định. Người tham gia được khuyến khích thực hiện các hoạt động theo cách tự nhiên nhất của họ: đi bộ trong hành lang, sân hoặc đường phố; đứng hoặc ngồi tại bất kỳ địa điểm nào trong sinh hoạt thường ngày. Cách tiếp cận này nhằm thu thập hành vi sử dụng điện thoại thực tế thay vì hành vi nhân tạo trong môi trường phòng thí nghiệm kiểm soát chặt, từ đó tăng khả năng tổng quát hóa của mô hình trong điều kiện triển khai thực tế.

Về phần cứng, bộ dữ liệu được thu thập trên 19 dòng thiết bị khác nhau từ năm hãng sản xuất phổ biến cùng một nhóm thiết bị khác, với phiên bản Android trải dài từ 8.1 đến 16. Bảng~\ref{tab:devices} tổng hợp thông tin về các thiết bị sử dụng.

\begin{table}[H]
    \centering
    \caption{Tổng hợp thiết bị tham gia thu thập dữ liệu}
    \label{tab:devices}
    \begin{tabular}{|l|c|c|}
    \hline
    \textbf{Hãng sản xuất} & \textbf{Số lượng máy} 
        & \textbf{Phiên bản Android} \\
    \hline
    Samsung  & 6 & 12, 13, 16 \\
    \hline
    Xiaomi   & 6 & 13, 14, 16 \\
    \hline
    OPPO     & 4 & 8.1, 12, 13 \\
    \hline
    Realme   & 1 & 15 \\
    \hline
    TECNO    & 1 & 15 \\
    \hline
    Khác     & 1 & 15 \\
    \hline
    \textbf{Tổng} & \textbf{19} & \textbf{8.1 - 16} \\
    \hline
    \end{tabular}
\end{table}

Sự đa dạng về phần cứng này là có chủ đích: các hãng khác nhau sử dụng chip cảm biến quán tính khác nhau, dẫn đến sự biến thiên về tần số lấy mẫu thực tế, độ nhiễu và đặc tính tín hiệu. Một mô hình được huấn luyện trên bộ dữ liệu đa thiết bị như vậy có tính thực tiễn cao hơn so với các công trình được thu thập trên một hoặc vài dòng máy đồng nhất, đây cũng là điểm khắc phục hạn chế được chỉ ra trong phân tích công trình liên quan tại Chương 2.

\section{Ứng dụng BioAuth Data Collection}
\label{sec:appthuthu}

Để thu thập dữ liệu một cách nhất quán trên nhiều người tham gia và nhiều thiết bị khác nhau, đề tài xây dựng ứng dụng Android chuyên dụng \textit{BioAuth Data Collection} (tên mã nội bộ DataCollectV2). Ứng dụng được phát triển bằng Kotlin 2.0.21, hỗ trợ từ Android 7.0 (API level 24) trở lên, và được tổ chức thành bốn màn hình chính theo một luồng tuyến tính: đăng ký thông tin, thu cảm biến quán tính, điền biểu mẫu tương tác và xuất dữ liệu.

\subsection{Màn hình đăng ký}

Trước khi bắt đầu thu thập, người tham gia khai báo thông tin cá nhân trên màn hình đăng ký, gồm tên, tuổi, giới tính và tay thuận; trong đó giới tính và tay thuận được chọn từ danh sách định sẵn nhằm tránh sai lệch khi nhập tự do. Thông tin thiết bị, gồm hãng sản xuất, model và phiên bản Android, được ứng dụng tự động ghi nhận từ hệ thống mà không yêu cầu người dùng tự nhập. Giao diện màn hình đăng ký được minh họa trong Hình~\ref{fig:dangky}.

\begin{figure}[H]
    \centering
    \includegraphics[width=0.28\textwidth,height=0.42\textheight,keepaspectratio]{Hinh_ve/Collectdata/screen_register.jpg}
    \caption{Giao diện màn hình đăng ký của ứng dụng BioAuth Data Collection}
    \label{fig:dangky}
\end{figure}

Như thể hiện trong Hình~\ref{fig:dangky}, màn hình đăng ký còn hiển thị một bản thông báo đồng thuận, nêu rõ ứng dụng chỉ thu dữ liệu cảm biến chuyển động và nhịp điệu thao tác, không lưu nội dung văn bản người dùng gõ, và dữ liệu chỉ phục vụ mục đích nghiên cứu. Người tham gia phải tích chọn ô đồng ý thì nút bắt đầu thu thập mới được kích hoạt, qua đó bảo đảm yêu cầu về đồng thuận có hiểu biết trong nghiên cứu trên người.

Khi hoàn tất, ứng dụng sinh một định danh \texttt{user\_id} duy nhất theo quy ước tiền tố \texttt{user\_} kèm Unix timestamp tại thời điểm đăng ký. Việc dùng mốc thời gian mili-giây đảm bảo không xảy ra xung đột định danh giữa các người tham gia, kể cả khi thu thập ở những thời điểm và trên những thiết bị khác nhau. Toàn bộ thông tin được lưu vào tệp \texttt{metadata.csv} theo định dạng hai cột là trường và giá trị, đặt trong thư mục riêng mang tên \texttt{user\_id} của người tham gia, làm cơ sở định danh cho toàn bộ dữ liệu cảm biến và tương tác thu được ở các bước sau.

\subsection{Màn hình thu cảm biến quán tính}

Màn hình thu cảm biến (\texttt{SensorCollectionActivity}) là thành phần trung tâm của ứng dụng, chịu trách nhiệm ghi nhận tín hiệu từ ba cảm biến quán tính là gia tốc kế, con quay hồi chuyển và cảm biến từ trường thông qua API \texttt{SensorManager} ở chế độ \texttt{SENSOR\_DELAY\_GAME}, tương ứng tần số lấy mẫu xấp xỉ 50 Hz và dao động tùy thiết bị. Người tham gia lần lượt thực hiện ba hoạt động đi bộ, đứng và ngồi với thời lượng mục tiêu mặc định 360 giây cho mỗi hoạt động, có thể điều chỉnh theo bước 30 giây trong khoảng từ 30 đến 1800 giây ngay trên giao diện.

Giao diện cung cấp phản hồi trực quan theo ba trạng thái của một phiên thu. Hình~\ref{fig:thucb} minh họa lần lượt trạng thái chưa thu, đang thu và đã hoàn tất.

\begin{figure}[H]
    \centering
    \begin{subfigure}[t]{0.3\textwidth}
        \centering
        \includegraphics[width=\textwidth,height=0.34\textheight,keepaspectratio]{Hinh_ve/Collectdata/screen_collect.jpg}
        \caption{Chưa thu}
        \label{fig:thucb_a}
    \end{subfigure}
    \hfill
    \begin{subfigure}[t]{0.3\textwidth}
        \centering
        \includegraphics[width=\textwidth,height=0.34\textheight,keepaspectratio]{Hinh_ve/Collectdata/screen_collecting.jpg}
        \caption{Đang thu}
        \label{fig:thucb_b}
    \end{subfigure}
    \hfill
    \begin{subfigure}[t]{0.3\textwidth}
        \centering
        \includegraphics[width=\textwidth,height=0.34\textheight,keepaspectratio]{Hinh_ve/Collectdata/screen_collected.jpg}
        \caption{Đã hoàn tất}
        \label{fig:thucb_c}
    \end{subfigure}
    \caption{Ba trạng thái của màn hình thu cảm biến: trước khi thu, trong khi thu và sau khi hoàn tất cả ba hoạt động}
    \label{fig:thucb}
\end{figure}

Ở trạng thái chưa thu (Hình~\ref{fig:thucb_a}), mỗi hoạt động hiển thị tiến độ 0\% cùng mốc thời lượng mục tiêu, và giá trị cảm biến thời gian thực ở mức 0. Khi một hoạt động được kích hoạt (Hình~\ref{fig:thucb_b}), ứng dụng hiển thị nhãn hoạt động đang thu, đồng hồ đếm thời gian, dải biểu đồ tín hiệu cập nhật liên tục và phần trăm tiến độ; một dòng nhắc nhở người dùng cầm máy tự nhiên thay vì giữ cố định nhằm thu được hành vi thực tế. Khi cả ba hoạt động đạt mục tiêu (Hình~\ref{fig:thucb_c}), tổng tiến độ đạt 100\% và ứng dụng cho phép chuyển sang màn hình biểu mẫu tương tác. Ô hiển thị giá trị thời gian thực chỉ trình bày gia tốc kế và con quay hồi chuyển cho gọn giao diện, song tín hiệu từ cảm biến từ trường vẫn được ghi nhận đầy đủ để tạo thành bộ dữ liệu chín kênh.

Để việc ghi nhận không bị gián đoạn khi người dùng chuyển ứng dụng hoặc tắt màn hình, quá trình thu được duy trì bởi một dịch vụ nền \texttt{SensorForegroundService} kiểu \texttt{START\_STICKY}. Mỗi lần thu của một hoạt động được lưu thành một tệp riêng theo quy ước \texttt{\{hoạt\_động\}\_att\{N\}.csv}, trong đó $N$ là số thứ tự lần thu tăng dần trong cùng một phiên; cơ chế này cho phép thu lại một hoạt động nhiều lần và đặt lại khi cần mà không ảnh hưởng đến dữ liệu của các hoạt động khác.

\subsection{Màn hình biểu mẫu tương tác}

Sau khi hoàn thành thu cảm biến, người tham gia điền biểu mẫu tương tác (\texttt{FormActivity}) được thiết kế gồm ba phần nhằm thu đồng thời ba loại hành vi tương tác màn hình. Phần một gồm 14 câu hỏi lựa chọn để thu động học chạm (tap dynamics); phần hai gồm 7 câu nhập văn bản tự do để thu động học gõ phím (keystroke dynamics); phần ba gồm 5 danh sách chọn nhiều kết hợp thao tác cuộn để thu động học cuộn (scroll dynamics). Hình~\ref{fig:form} minh họa phần đầu của biểu mẫu.

\begin{figure}[H]
    \centering
    \includegraphics[width=0.28\textwidth,height=0.42\textheight,keepaspectratio]{Hinh_ve/Collectdata/screen_form.jpg}
    \caption{Giao diện màn hình biểu mẫu tương tác, phần câu hỏi lựa chọn}
    \label{fig:form}
\end{figure}

Các câu hỏi được thiết kế gần gũi với thói quen sinh hoạt thường ngày, như thể hiện trong Hình~\ref{fig:form}, nhằm khuyến khích người tham gia thao tác tự nhiên. Điểm cần lưu ý là nội dung câu trả lời không phải là đối tượng quan tâm của đề tài; thay vào đó, ứng dụng ghi lại các đặc trưng hành vi của thao tác. Với mỗi lần chạm, ứng dụng ghi nhận các sự kiện chạm xuống và nhấc tay cùng tọa độ, dấu thời gian và thời gian giữ; với mỗi lần gõ phím, ứng dụng ghi nhận thời điểm gõ để tính các đặc trưng nhịp điệu; với thao tác cuộn, ứng dụng ghi nhận quỹ đạo và vận tốc cuộn. Cách thu thông qua các tác vụ chuẩn hóa như vậy bảo đảm tính nhất quán giữa các người tham gia, đồng thời không cần đến nội dung văn bản người dùng nhập, phù hợp với cam kết về quyền riêng tư đã nêu ở màn hình đăng ký.

Mỗi phiên cho phép điền biểu mẫu tối đa hai lượt nhằm tăng khối lượng dữ liệu tương tác. Ba loại sự kiện được lưu thành ba tệp riêng theo quy ước \texttt{tap\_r\{N\}.csv}, \texttt{keystroke\_r\{N\}.csv} và \texttt{scroll\_r\{N\}.csv}, trong đó $N$ là số thứ tự lượt điền trong phiên.

\subsection{Màn hình xuất dữ liệu và cấu trúc dữ liệu đầu ra}

Sau khi hoàn thành thu cảm biến và biểu mẫu, ứng dụng hiển thị màn hình tổng kết phiên thu (\texttt{UploadActivity}), như minh họa trong Hình~\ref{fig:xuatdulieu}. Màn hình này hiển thị định danh \texttt{user\_id}, phiên hiện tại và bảng thống kê dữ liệu đã thu được, gồm số tệp CSV tín hiệu quán tính cùng số mẫu sự kiện chạm, gõ phím và cuộn, kèm dung lượng ước tính của phiên.

\begin{figure}[H]
    \centering
    \includegraphics[width=0.28\textwidth,height=0.42\textheight,keepaspectratio]{Hinh_ve/Collectdata/screen_finshed.jpg}
    \caption{Màn hình tổng kết và xuất dữ liệu của một phiên thu}
    \label{fig:xuatdulieu}
\end{figure}

Như thể hiện trong Hình~\ref{fig:xuatdulieu}, ứng dụng kiểm tra tính đầy đủ của dữ liệu trước khi cho phép xuất, cảnh báo nếu một phiên chưa có đủ tệp tín hiệu quán tính cho cả ba hoạt động. Khi xuất, toàn bộ dữ liệu của phiên được nén thành một tệp ZIP đặt tên theo \texttt{user\_id} kèm dấu thời gian, sau đó được chia sẻ thông qua cơ chế chia sẻ tệp chuẩn của Android tới các ứng dụng như Gmail hoặc Google Drive. Cách này cho phép thu thập phân tán trên nhiều thiết bị và địa điểm khác nhau mà không cần máy chủ trung gian.

Toàn bộ dữ liệu của một người tham gia được tổ chức theo cấu trúc phân cấp \texttt{user\_\{id\}/session\_\{N\}/}, trong đó mỗi thư mục phiên chứa đầy đủ các tệp tín hiệu quán tính và tương tác màn hình của phiên đó. Sau khi hoàn tất một phiên, ứng dụng tự động tăng chỉ số phiên để chuẩn bị cho lần thu tiếp theo, nhờ đó mỗi người tham gia có thể đóng góp nhiều phiên độc lập theo thời gian. Cấu trúc thư mục nhất quán này tạo thuận lợi cho bước tiền xử lý và gán nhãn dữ liệu được trình bày ở phần tiếp theo.

\section{Tiền xử lý dữ liệu}
\label{sec:tienxuly}

Dữ liệu cảm biến thô thu được từ ứng dụng không thể đưa trực tiếp vào mô hình mà cần qua một quy trình tiền xử lý nhằm loại bỏ dữ liệu kém chất lượng, phân đoạn thành các đơn vị đầu vào có kích thước cố định và chuẩn hóa để giảm ảnh hưởng của khác biệt phần cứng. Phần này mô tả quy trình tiền xử lý dữ liệu quán tính chín kênh, gồm gia tốc kế, con quay hồi chuyển và cảm biến từ trường theo ba trục.

\subsection{Kiểm tra chất lượng và phân đoạn tín hiệu}

Mỗi tệp tín hiệu được kiểm tra tính liên tục theo dấu thời gian. Do tần số lấy mẫu thực tế dao động giữa các thiết bị, một số mẫu có thể bị mất hoặc trễ; những vị trí có khoảng trống thời gian vượt quá ngưỡng tương ứng năm chu kỳ lấy mẫu được xem là điểm gián đoạn và chuỗi tín hiệu được tách thành các đoạn liên tục tại đó. Các đoạn ngắn hơn độ dài một cửa sổ sẽ bị loại bỏ, và thời lượng hợp lệ của mỗi hoạt động chỉ được tính trên phần tín hiệu liên tục. Quy trình kiểm tra còn rà soát những phiên có thống kê bất thường so với phần còn lại của cùng người dùng dựa trên độ lệch tuyệt đối trung vị, nhằm phát hiện và xử lý sớm các phiên thu lỗi trước khi đưa vào huấn luyện.

\subsection{Phân cửa sổ và chuẩn hóa}

Sau khi phân đoạn, mỗi đoạn tín hiệu liên tục được cắt thành các cửa sổ có độ dài cố định 200 mẫu, tương ứng khoảng bốn giây ở tần số 50 Hz, bằng kỹ thuật cửa sổ trượt với bước trượt 20 mẫu. Độ dài bốn giây được chọn để mỗi cửa sổ chứa được vài chu kỳ của các hoạt động có tính tuần hoàn như đi bộ, đủ để mô hình quan sát một mẫu hành vi hoàn chỉnh, đồng thời không quá dài khiến độ trễ ra quyết định khi vận hành tăng cao. Bước trượt 20 mẫu tạo mức chồng lấn lớn giữa các cửa sổ liên tiếp; việc cho các cửa sổ chồng lấn giúp tăng đáng kể số lượng mẫu huấn luyện và giúp mô hình quan sát cùng một hành vi ở nhiều vị trí thời gian khác nhau, qua đó tăng tính bền vững với độ lệch pha của tín hiệu. Mỗi cửa sổ là một ma trận kích thước $200 \times 9$ tương ứng 200 bước thời gian trên chín kênh cảm biến. Hình~\ref{fig:cua_so_truot} minh họa trực quan cơ chế cửa sổ trượt: các cửa sổ liên tiếp chỉ dịch đi 20 mẫu nên chồng lấn phần lớn lên nhau, qua đó vừa tăng số lượng mẫu huấn luyện, vừa cho mô hình quan sát cùng một hành vi ở nhiều vị trí thời gian khác nhau.

\begin{figure}[H]
  \centering
  \includegraphics[width=0.85\textwidth]{Hinh_ve/cua_so_truot.png}
  \caption{Minh họa kỹ thuật cửa sổ trượt trên tín hiệu gia tốc kế}
  \label{fig:cua_so_truot}
\end{figure}

Mỗi cửa sổ được chuẩn hóa độc lập theo phương pháp z-score dọc theo trục thời gian, tức là trên từng kênh, giá trị được trừ đi trung bình và chia cho độ lệch chuẩn tính trong phạm vi cửa sổ đó. Cụ thể, với kênh $c$ và bước thời gian $t$ trong một cửa sổ, giá trị sau chuẩn hóa được tính theo

\begin{equation}
    \tilde{x}_{t,c} = \frac{x_{t,c} - \mu_c}{\sigma_c + \epsilon},
    \label{eq:zscore}
\end{equation}

trong đó $x_{t,c}$ là giá trị gốc, $\mu_c$ và $\sigma_c$ lần lượt là trung bình và độ lệch chuẩn của kênh $c$ tính trên toàn bộ cửa sổ, và $\epsilon$ là một hằng số dương rất nhỏ thêm vào mẫu số để tránh chia cho không. Cách chuẩn hóa theo từng cửa sổ này đưa tín hiệu của mọi thiết bị về cùng một thang đo, qua đó giảm bớt ảnh hưởng của khác biệt về thang đo và độ lệch nền giữa các chip cảm biến khác nhau, một yếu tố quan trọng đối với bộ dữ liệu đa thiết bị của đề tài.

Một điểm cần lưu ý là phép chuẩn hóa này cũng loại bỏ một phần thông tin của tín hiệu. Việc trừ trung bình theo từng kênh gột đi thành phần một chiều trong mỗi cửa sổ - chủ yếu là vector trọng lực, vốn mã hóa hướng nghiêng và tư thế cầm máy - còn phép chia cho độ lệch chuẩn loại bỏ thông tin về biên độ tuyệt đối trên từng trục. Đây là một đánh đổi có chủ đích: hướng và tư thế cầm máy bản thân chúng có thể mang thông tin phân biệt người dùng, nhưng cũng biến thiên mạnh theo ngữ cảnh và phụ thuộc vào đặc tính của từng cảm biến, nên việc gột bỏ giúp tín hiệu ổn định hơn giữa các dòng thiết bị - đổi lại, hệ thống không khai thác đặc trưng hướng và tư thế. Việc giữ lại có chọn lọc thông tin này, chẳng hạn tách riêng thành phần trọng lực thay vì trừ trung bình toàn cục, được xem là một hướng cải thiện trong tương lai.

Hình~\ref{fig:zscore_scale} minh họa tác dụng rõ rệt nhất của chuẩn hóa z-score lên thang đo của chín kênh cảm biến. Trước chuẩn hóa (Hình~\ref{fig:scale_truoc}), giá trị thô của ba nhóm cảm biến chênh lệch rất lớn về biên độ: cảm biến từ trường trải tới vài chục đơn vị trong khi con quay hồi chuyển gần như bằng không, khiến những kênh có biên độ lớn dễ lấn át khi mô hình học. Sau chuẩn hóa theo từng cửa sổ (Hình~\ref{fig:scale_sau}), mọi kênh được đưa về cùng một dải phân bố quanh trung bình 0 với độ lệch chuẩn đơn vị, bảo đảm các kênh đóng góp cân bằng.

\begin{figure}[H]
  \centering
  \begin{subfigure}{0.49\textwidth}\centering
    \includegraphics[width=\linewidth]{Hinh_ve/3a_scale_truoc.png}
    \caption{Trước chuẩn hóa}\label{fig:scale_truoc}
  \end{subfigure}\hfill
  \begin{subfigure}{0.49\textwidth}\centering
    \includegraphics[width=\linewidth]{Hinh_ve/3b_scale_sau.png}
    \caption{Sau chuẩn hóa}\label{fig:scale_sau}
  \end{subfigure}
  \caption{Phân bố giá trị 9 kênh cảm biến trước và sau chuẩn hóa z-score}
  \label{fig:zscore_scale}
\end{figure}

Bên cạnh việc thống nhất thang đo giữa các kênh, chuẩn hóa theo từng cửa sổ còn giúp căn chỉnh tín hiệu giữa những người dùng và thiết bị khác nhau. Hình~\ref{fig:zscore_user} so sánh cùng kênh gia tốc khi đi bộ của ba người dùng: tín hiệu thô lệch nhau về độ lệch nền và biên độ do cách cầm máy và đặc tính cảm biến khác nhau (Hình~\ref{fig:zscore_user_truoc}), nhưng sau chuẩn hóa thì được căn về cùng một dải quanh giá trị 0 trong khi vẫn giữ nguyên hình dạng dao động đặc trưng của hành vi (Hình~\ref{fig:zscore_user_sau}).

\begin{figure}[H]
  \centering
  \begin{subfigure}{0.49\textwidth}\centering
    \includegraphics[width=\linewidth]{Hinh_ve/2a_zscore_truoc.png}
    \caption{Trước chuẩn hóa}\label{fig:zscore_user_truoc}
  \end{subfigure}\hfill
  \begin{subfigure}{0.49\textwidth}\centering
    \includegraphics[width=\linewidth]{Hinh_ve/2b_zscore_sau.png}
    \caption{Sau chuẩn hóa z-score}\label{fig:zscore_user_sau}
  \end{subfigure}
  \caption{Tín hiệu gia tốc \texttt{acc\_x} của ba người dùng trước và sau chuẩn hóa z-score}
  \label{fig:zscore_user}
\end{figure}

Để minh họa sự khác biệt về bản chất tín hiệu giữa các hoạt động, Hình~\ref{fig:tin_hieu_tho} biểu diễn tín hiệu gia tốc kế thô trong mười giây đầu của một phiên ở hai hoạt động tiêu biểu. Khi đi bộ (Hình~\ref{fig:tin_hieu_di_bo}), ba trục gia tốc dao động tuần hoàn rõ rệt với chu kỳ tương ứng nhịp bước chân khoảng hai bước mỗi giây và biên độ lớn, phản ánh chuyển động có cấu trúc của cơ thể. Ngược lại, khi ngồi (Hình~\ref{fig:tin_hieu_ngoi}), tín hiệu gần như là đường thẳng quanh giá trị tĩnh, biến thiên rất nhỏ. Sự tương phản này cho thấy hoạt động động mang nhiều thông tin phân biệt người dùng hơn hoạt động tĩnh, và là cơ sở trực quan cho việc khảo sát riêng cấu hình chỉ đi bộ bên cạnh cấu hình toàn bộ hoạt động. Hình sử dụng tín hiệu thô trước khi chuẩn hóa để giữ nguyên đơn vị vật lý và thành phần trọng lực, nhờ đó thể hiện rõ bản chất chuyển động.

\begin{figure}[H]
  \centering
  \begin{subfigure}{0.49\textwidth}\centering
    \includegraphics[width=\linewidth]{tin_hieu_di_bo.png}
    \caption{Đi bộ}\label{fig:tin_hieu_di_bo}
  \end{subfigure}\hfill
  \begin{subfigure}{0.49\textwidth}\centering
    \includegraphics[width=\linewidth]{tin_hieu_ngoi.png}
    \caption{Ngồi}\label{fig:tin_hieu_ngoi}
  \end{subfigure}
  \caption{Tín hiệu gia tốc kế thô (trước chuẩn hóa) của một cửa sổ bốn giây}
  \label{fig:tin_hieu_tho}
\end{figure}

\subsection{Tổ chức dữ liệu theo ngữ cảnh hoạt động}

Như đã đề cập ở Chương 2, đề tài khai thác ngữ cảnh hoạt động bằng cách huấn luyện và đánh giá theo từng cấu hình ngữ cảnh thay vì xây dựng một tầng phân loại hoạt động trực tuyến. Do tín hiệu ở hai hoạt động đứng và ngồi gần như không biến thiên và khó phân biệt, bước tiền xử lý xuất ra hai tập dữ liệu song song: một tập chỉ gồm dữ liệu đi bộ và một tập gồm cả ba hoạt động. Hai tập này được sử dụng để khảo sát ảnh hưởng của ngữ cảnh hoạt động tới hiệu năng xác thực, và kết quả so sánh được trình bày ở Chương 5. Mỗi cửa sổ sau khi chuẩn hóa được gắn nhãn theo định danh người dùng, tạo thành các tensor đầu vào sẵn sàng cho giai đoạn huấn luyện mô hình.

Cần lưu ý rằng dữ liệu quán tính được đưa trực tiếp vào mô hình dưới dạng tín hiệu thô đã chuẩn hóa, không qua bước trích chọn đặc trưng thủ công, nhằm để mạng học sâu tự học các đặc trưng phân biệt từ dữ liệu. Ngược lại, dữ liệu tương tác màn hình cần được chuyển thành các đặc trưng động học trước khi phân loại; quá trình trích chọn đặc trưng cho dữ liệu này được trình bày cùng với nhánh mô hình tương ứng trong Chương 4.

\subsection{Phân tích dữ liệu sau tiền xử lý}

Sau khi hoàn tất tiền xử lý, bộ dữ liệu được khảo sát để hiểu rõ quy mô và đặc điểm trước khi huấn luyện. Hình~\ref{fig:phanbo} trình bày phân bố số cửa sổ thu được, xét theo từng người dùng và theo cấu hình ngữ cảnh. Số cửa sổ của mỗi người dùng (Hình~\ref{fig:phanbo_nguoi}) không đồng đều, dao động từ khoảng ba nghìn đến hơn tám nghìn cửa sổ tùy theo số phiên và thời lượng mỗi người tham gia thu thập. Xét theo cấu hình ngữ cảnh (Hình~\ref{fig:phanbo_cauhinh}), cấu hình toàn bộ hoạt động gồm 136.746 cửa sổ từ 26 người dùng, trong khi cấu hình chỉ đi bộ gồm 45.899 cửa sổ; tỉ lệ khoảng một phần ba phản ánh việc cấu hình toàn bộ gộp cả ba hoạt động còn cấu hình đi bộ chỉ giữ một hoạt động. Sự mất cân bằng số cửa sổ giữa các người dùng là lý do đề tài ưu tiên các chỉ số đánh giá không phụ thuộc phân bố lớp như tỉ lệ lỗi cân bằng, thay vì chỉ dựa vào độ chính xác vốn dễ bị chi phối bởi những người dùng có nhiều cửa sổ.

\begin{figure}[H]
  \centering
  \begin{subfigure}{0.62\textwidth}\centering
    \includegraphics[width=\linewidth]{Hinh_ve/phanbo_theo_nguoi.png}
    \caption{Theo người dùng}\label{fig:phanbo_nguoi}
  \end{subfigure}\hfill
  \begin{subfigure}{0.36\textwidth}\centering
    \includegraphics[width=\linewidth]{Hinh_ve/phanbo_theo_cauhinh.png}
    \caption{Theo cấu hình}\label{fig:phanbo_cauhinh}
  \end{subfigure}
  \caption{Phân bố số cửa sổ sau tiền xử lý}
  \label{fig:phanbo}
\end{figure}

Chương này đã trình bày toàn bộ quá trình xây dựng dữ liệu cho hệ thống, từ thiết kế thu thập, ứng dụng thu dữ liệu chuyên dụng đến quy trình tiền xử lý tín hiệu quán tính. Điểm nhấn trong thiết kế là thu hành vi trong điều kiện sử dụng tự nhiên trên nhiều dòng thiết bị, kết hợp quy trình tiền xử lý gồm kiểm tra chất lượng, phân đoạn, phân cửa sổ trượt và chuẩn hóa theo từng cửa sổ để tạo ra bộ dữ liệu nhất quán và có khả năng tổng quát hóa cao. Kết quả của chương là các tensor tín hiệu quán tính đã chuẩn hóa, tổ chức theo hai cấu hình ngữ cảnh và gắn nhãn theo người dùng. Trên cơ sở dữ liệu này, Chương 4 trình bày quá trình xây dựng mô hình, bao gồm so sánh các kiến trúc backbone cho dữ liệu quán tính, trích xuất đặc trưng và phân loại dữ liệu tương tác màn hình, cơ chế dung hợp và hàm quyết định nhận dạng tập mở, cùng việc xuất mô hình sang định dạng TensorFlow Lite.

\end{document}